\section{Results in four dimensions}
\label{sec:results4}

\subsection{Exhaustive sweeps}

Horv\'{a}th and L\'{a}ngi extended Property Z to all dimensions, proved that
maximisers are simplicial, and solved $n = d+2$, together with $n = d+3$ for
odd $d$ and among noncyclic polytopes for even $d$; they also proved that the
maximal-volume cross-polytope inscribed in $\mathbb{S}^{d-1}$ is the regular
one \cite{horvath-langi-2016}. Because a maximiser is simplicial and
Firsching's enumeration is complete for $n \leq 8$, a per-type sweep addresses
numerically the cross-class question those results leave open \artref{cases}.

\begin{center}
\begin{tabular}{lrrrr}
  \toprule
  $n$ & 5 & 6 & 7 & 8 \\
  \midrule
  types  & 1 & 2 & 5 & 37 \\
  winner & $25\sqrt{5}/384$ & $9/32$ & $\sqrt{3}/4$ & $2/3$ \\
  level  & E1 & E3 & E3 & E3 \\
  \bottomrule
\end{tabular}
\end{center}

At $n = 5$ the enumeration holds a single type, so as at $d = 3$, $n = 4, 5$
there is no rival to exclude and the case stays at E1. Two of the others
deserve comment.

At $n = 7$ the winner over all five types is $\sqrt{3}/4 = 0.433012701892$, a
triangle and a square in orthogonal planes. The cyclic type $C_{4}(7)$, degree
sequence $[6]^{7}$, maximises at $0.389054556345$, agreeing with the
Horv\'{a}th--L\'{a}ngi closed form $(49/192)(\cos(\pi/7) + \cos(2\pi/7))$ to
$6 \times 10^{-17}$. Because this is the maximum found over the whole cyclic
type rather than the value at one configuration, it is direct numerical
evidence for their Problem 1, that the maximal-volume cyclic polytopes in
$\mathcal{P}_{d}(d+3)$ are the congruent copies of $C_{d}(d+3)$, and for their
Problem 2, that for even $d$ the overall maximiser is not cyclic.

At $n = 8$ all $37$ types are swept and the regular cross-polytope wins at
$2/3$, ahead of the best rival in convex position, $0.533173$
(\texttt{firsching-8-5}), by $20.0\%$ of its own value. The triage strengthens
this: every rival's \emph{signed} maximum also falls below $2/3$, the largest
being $0.601260$ (\texttt{firsching-8-22}) with a $9.81\%$ margin, so the case
reaches E3. Horv\'{a}th and L\'{a}ngi prove optimality here only within the
cross-polytope class, whereas the sweep ranges over every combinatorial type,
so it covers a wider class of polytopes, though only numerically.

\subsection{The orthogonal polygon-pair family}

The winners at $n = 6, 7, 8$ share one description: $n = 6$ is two triangles,
$n = 7$ a triangle and a square, and $n = 8$ two squares, the $16$-cell, each
pair lying in orthogonal $2$-planes. Since the polygons lie in complementary
orthogonal subspaces through the origin their hull is a free sum, so with
$\operatorname{area}(P_{m}) = (m/2)\sin(2\pi/m)$,
\begin{equation}
  \label{eq:polygonpair}
  \vol_{4}(m,k)
  = \frac{2!\,2!}{4!}\operatorname{area}(P_{m})\operatorname{area}(P_{k})
  = \frac{mk}{24}\sin\frac{2\pi}{m}\sin\frac{2\pi}{k},
  \qquad m + k = n .
\end{equation}
This requires $m, k \geq 3$: a $2$-gon is a segment and spans one dimension
rather than two, which is why $n = 5$ falls outside the family and is the
regular simplex instead.

\begin{table}[ht]
  \caption{The polygon-pair family. ``Status'' records the published status;
    \emph{class-proven} means optimality was proved only inside a restricted
    class, noncyclic types at $n = 7$ and cross-polytopes at $n = 8$.}
  \label{tab:pair}
  \begin{tabular}{rlllr}
    \toprule
    $n$ & split & exact volume & numeric & status \\
    \midrule
    6  & $(3,3)$ & $9/32$ & 0.281250000000 & proven \\
    7  & $(3,4)$ & $\sqrt{3}/4$ & 0.433012701892 & class-proven \\
    8  & $(4,4)$ & $2/3$ & 0.666666666667 & class-proven \\
    \bottomrule
  \end{tabular}
\end{table}

Each row of Table~\ref{tab:pair} is the best split of its $n$ and equals the
sweep winner above. Free search \artref{free} matches every row to $10^{-15}$,
and each value is verified against the explicitly constructed polytope
\artref{families} \artref{testtables}.

The family does not stay optimal. The type-free search \artref{free} matches
the best split at $n = 9$ and $n = 10$ to $10^{-15}$, but at $n = 11, 12, 13$
it finds inscribed polytopes with all $n$ points as vertices whose volumes
$1.046529$, $1.171382$ and $1.303558$ exceed the best split by $1.65\%$,
$4.12\%$ and $10.01\%$ \artref{testtables}. Each is an explicit configuration,
so the family is not optimal for these $n$; at $n = 9, 10$, where the
enumeration of types is no longer complete, its optimality is neither shown nor
refuted.
